\documentclass[11pt]{article}

\usepackage[margin=1in]{geometry}
\usepackage{amsmath,amssymb,amsthm}
\usepackage{hyperref}

\hypersetup{hidelinks}

\newtheorem{theorem}{Theorem}[section]
\newtheorem{lemma}[theorem]{Lemma}
\newtheorem{proposition}[theorem]{Proposition}
\newtheorem{conjecture}[theorem]{Conjecture}
\newtheorem{corollary}[theorem]{Corollary}
\newtheorem{remark}[theorem]{Remark}

\newcommand{\N}{\mathbb{N}}
\newcommand{\Z}{\mathbb{Z}}
\newcommand{\calS}{\mathcal{S}}
\newcommand{\rfour}{r_4}
\newcommand{\rthree}{r_3}

\title{A Partial Van Corput Reduction for Four-Term Progression-Free Sets}
\author{Partial-results manuscript}
\date{June 2026}

\begin{document}
\maketitle

\begin{abstract}
Let \(r_k(N)\) denote the maximum size of a subset of \(\{1,\ldots,N\}\) with no
non-trivial \(k\)-term arithmetic progression.  We record a partial van
Corput-type reduction from the \(k=4\) problem to the \(k=3\) problem.  For a
4-term-progression-free set \(A\subseteq\{1,\ldots,N\}\), write
\[
  A_d=\{x\in A:x+d\in A\}.
\]
We prove that there is a popular difference \(d\) satisfying
\[
  |A_d|\geq |A|^2/(4N),
\]
and that such a fiber avoids 3-term progressions of step \(d\) and \(2d\).
The missing ingredient for a full reduction is the assertion that some popular
fiber is free of all 3-term progressions.  We isolate this as Conjecture 2.
Assuming Conjecture 2, one obtains
\[
  r_4(N)\leq 2\sqrt{N r_3(N)}.
\]
Combining this conditional reduction with Raghavan's 2026 bound for \(r_3(N)\)
would give a conditional quasi-polynomial upper bound for \(r_4(N)\).
We also describe a \(2\times3\) grid reformulation of the gap, Lean 4
formalization status, and computational evidence supporting the conjecture.
\end{abstract}

\tableofcontents

\section{Introduction}

For \(k\geq 2\), let \(r_k(N)\) be the largest size of a subset
\[
  A\subseteq \{1,\ldots,N\}
\]
containing no non-trivial \(k\)-term arithmetic progression.  Thus \(A\) is
\(k\)-AP-free if it contains no set
\[
  \{a,a+d,\ldots,a+(k-1)d\}
\]
with \(d>0\).

Szemerédi's theorem gives \(r_k(N)=o(N)\) for each fixed \(k\geq 3\), but the
quantitative behavior remains difficult.  The \(k=3\) case has seen major
recent progress.  The input used in this note is Raghavan's 2026 bound:
\[
  r_3(N)
  \leq
  C_1 N
  \exp\!\left(
    -c_1\frac{(\log N)^{1/6}}{\log\log N}
  \right)
\]
for absolute constants \(C_1,c_1>0\) and all sufficiently large \(N\).

For \(k=4\), the best unconditional published upper bounds are much weaker than
this quasi-polynomial scale.  The purpose of this note is not to claim an
unconditional improvement, but to isolate a precise combinatorial condition
under which Raghavan's \(r_3(N)\) bound transfers to \(r_4(N)\).  The transfer
rests on popular difference fibers.

\section{Setup and Difference Fibers}

Let
\[
  A\subseteq\{1,\ldots,N\},\qquad |A|=M.
\]
For \(1\leq d\leq N-1\), define the difference fiber
\[
  A_d=\{x\in A:x+d\in A\}.
\]
Equivalently, \(A_d=A\cap(A-d)\), written inside \(\{1,\ldots,N\}\).

The first observation is a pigeonhole count.  It is elementary, but it is the
starting point of the whole reduction.

\begin{lemma}[Popular difference]
\label{lem:popular}
Let \(A\subseteq\{1,\ldots,N\}\) have size \(M\geq 2\).  Then there exists
\(d\in\{1,\ldots,N-1\}\) such that
\[
  |A_d|\geq \frac{M^2}{4N}.
\]
\end{lemma}

\begin{proof}
Count ordered pairs \((a,b)\in A^2\) with \(a<b\).  Every such pair determines a
unique positive difference \(d=b-a\), and contributes one element \(a\) to
\(A_d\).  Therefore
\[
  \sum_{d=1}^{N-1}|A_d|
  =
  \binom{M}{2}
  =
  \frac{M(M-1)}{2}.
\]
Since \(M\geq 2\),
\[
  \frac{M(M-1)}{2}\geq \frac{M^2}{4}.
\]
There are at most \(N-1<N\) possible positive differences.  Hence some \(d\)
satisfies
\[
  |A_d|
  \geq
  \frac{M^2/4}{N}
  =
  \frac{M^2}{4N}.
\]
\end{proof}

We call a difference \(d\) \emph{popular} if
\[
  |A_d|\geq \frac{M^2}{4N}.
\]
Let
\[
  \calS(A)=\left\{d\in\{1,\ldots,N-1\}: |A_d|\geq \frac{M^2}{4N}\right\}
\]
be the set of popular differences.

\section{What 4-AP-Freeness Forces}

Suppose now that \(A\) is 4-AP-free.  The key unconditional structural fact is
that fibers \(A_d\) avoid certain specific 3-term progressions.

\begin{lemma}[Step-\(d\) freeness]
\label{lem:stepd}
Let \(A\subseteq\{1,\ldots,N\}\) be 4-AP-free and let \(d>0\).  Then \(A_d\)
contains no 3-term arithmetic progression with common difference \(d\).  In
other words, there is no \(x\) such that
\[
  x,\ x+d,\ x+2d\in A_d.
\]
\end{lemma}

\begin{proof}
Assume for contradiction that \(x,x+d,x+2d\in A_d\).  By definition of \(A_d\),
\[
  x\in A_d \quad\Rightarrow\quad x\in A,\ x+d\in A,
\]
\[
  x+d\in A_d \quad\Rightarrow\quad x+d\in A,\ x+2d\in A,
\]
and
\[
  x+2d\in A_d \quad\Rightarrow\quad x+2d\in A,\ x+3d\in A.
\]
Thus
\[
  x,\ x+d,\ x+2d,\ x+3d\in A,
\]
which is a non-trivial 4-term arithmetic progression in \(A\).  This contradicts
the assumption that \(A\) is 4-AP-free.
\end{proof}

\begin{lemma}[Step-\(2d\) freeness]
\label{lem:step2d}
Let \(A\subseteq\{1,\ldots,N\}\) be 4-AP-free and let \(d>0\).  Then \(A_d\)
contains no 3-term arithmetic progression with common difference \(2d\).  In
other words, there is no \(x\) such that
\[
  x,\ x+2d,\ x+4d\in A_d.
\]
\end{lemma}

\begin{proof}
Assume for contradiction that \(x,x+2d,x+4d\in A_d\).  From \(x\in A_d\), we get
\[
  x\in A,\qquad x+d\in A.
\]
From \(x+2d\in A_d\), we get
\[
  x+2d\in A,\qquad x+3d\in A.
\]
Therefore
\[
  x,\ x+d,\ x+2d,\ x+3d\in A,
\]
again a 4-term arithmetic progression in \(A\).  This is impossible.
\end{proof}

Combining Lemmas \ref{lem:popular}, \ref{lem:stepd}, and \ref{lem:step2d} gives
the main unconditional partial result.

\begin{proposition}[Partial van Corput fiber structure]
\label{prop:partial}
Let \(A\subseteq\{1,\ldots,N\}\) be 4-AP-free and let \(|A|=M\geq 2\).  Then
there exists \(d\in\{1,\ldots,N-1\}\) such that:
\begin{enumerate}
  \item \( |A_d|\geq M^2/(4N)\);
  \item \(A_d\) contains no 3-term arithmetic progression with common difference \(d\);
  \item \(A_d\) contains no 3-term arithmetic progression with common difference \(2d\).
\end{enumerate}
\end{proposition}

\begin{proof}
Choose \(d\) by Lemma \ref{lem:popular}.  Since \(A\) is 4-AP-free, Lemma
\ref{lem:stepd} gives property (2), and Lemma \ref{lem:step2d} gives property
(3).
\end{proof}

\section{The Exact Gap}

Proposition \ref{prop:partial} is not enough to apply \(r_3(N)\), because
\(r_3(N)\) bounds sets with no 3-term arithmetic progressions of any positive
common difference.  The proposition only rules out common differences \(d\) and
\(2d\) inside \(A_d\).

This distinction is essential.

\begin{remark}[A warning example]
Let
\[
  A=\{1,2,3,5,6,8,9\}\subseteq\{1,\ldots,9\}.
\]
This set is 4-AP-free.  For \(d=1\),
\[
  A_1=\{1,2,5,8\}.
\]
But \(A_1\) contains the 3-term progression
\[
  2,\ 5,\ 8,
\]
whose common difference is \(3\).  Thus even a natural fiber of a 4-AP-free set
need not be fully 3-AP-free.
\end{remark}

The missing assertion is the following.

\begin{conjecture}[Conjecture 2: popular fully 3-AP-free fiber]
\label{conj:c2}
Let \(A\subseteq\{1,\ldots,N\}\) be 4-AP-free, with \(|A|=M\geq 2\).  Then
there exists \(d\in\calS(A)\) such that \(A_d\) is fully 3-AP-free.
\end{conjecture}

Equivalently, among the popular differences, at least one fiber has no
non-trivial 3-term progression with any positive common difference.

\section{Conditional Reduction from \(r_4\) to \(r_3\)}

Conjecture \ref{conj:c2} gives the desired reduction immediately.

\begin{theorem}[Conditional van Corput reduction]
\label{thm:conditional_vc}
Assume Conjecture \ref{conj:c2}.  Then for every \(N\),
\[
  r_4(N)\leq 2\sqrt{N r_3(N)}.
\]
\end{theorem}

\begin{proof}
Let \(A\subseteq\{1,\ldots,N\}\) be 4-AP-free and set \(M=|A|\).  If \(M<2\), the
claim is trivial.  Assume \(M\geq 2\).

By Conjecture \ref{conj:c2}, there exists a popular \(d\) such that \(A_d\) is
fully 3-AP-free.  Since \(d\) is popular,
\[
  |A_d|\geq \frac{M^2}{4N}.
\]
Since \(A_d\subseteq \{1,\ldots,N\}\) and is 3-AP-free,
\[
  |A_d|\leq r_3(N).
\]
Therefore
\[
  \frac{M^2}{4N}\leq r_3(N),
\]
so
\[
  M\leq 2\sqrt{N r_3(N)}.
\]
Taking \(A\) to have maximum possible size \(M=r_4(N)\) gives the result.
\end{proof}

\begin{corollary}[Conditional quasi-polynomial \(r_4\) bound]
\label{cor:conditional_r4}
Assume Conjecture \ref{conj:c2}.  Assume also Raghavan's bound
\[
  r_3(N)
  \leq
  C_1 N
  \exp\!\left(
    -c_1\frac{(\log N)^{1/6}}{\log\log N}
  \right)
\]
for all sufficiently large \(N\).  Then there exist constants \(C,c>0\) such
that for all sufficiently large \(N\),
\[
  r_4(N)
  \leq
  C N
  \exp\!\left(
    -c\frac{(\log N)^{1/6}}{\log\log N}
  \right).
\]
\end{corollary}

\begin{proof}
By Theorem \ref{thm:conditional_vc},
\[
  r_4(N)\leq 2\sqrt{N r_3(N)}.
\]
Using Raghavan's bound,
\[
  r_4(N)
  \leq
  2\sqrt{
    N\cdot
    C_1 N
    \exp\!\left(
      -c_1\frac{(\log N)^{1/6}}{\log\log N}
    \right)
  }.
\]
Thus
\[
  r_4(N)
  \leq
  2\sqrt{C_1}\,N
  \exp\!\left(
    -\frac{c_1}{2}\frac{(\log N)^{1/6}}{\log\log N}
  \right).
\]
Set \(C=2\sqrt{C_1}\) and \(c=c_1/2\).
\end{proof}

\begin{remark}
The square root changes the constant \(c_1\) to \(c_1/2\), but it does not change
the exponent \(1/6\) on \(\log N\).
\end{remark}

\section{A \(2\times3\) Grid Reformulation}

The remaining obstruction can be expressed as a grid-avoidance condition.  A
3-term progression in \(A_d\) with step \(e\) means
\[
  x,\ x+e,\ x+2e\in A_d.
\]
By definition of \(A_d\), this is equivalent to the six points
\[
  x,\ x+e,\ x+2e,\ x+d,\ x+e+d,\ x+2e+d
\]
all lying in \(A\).  These six points form a \(2\times 3\) arithmetic grid:
\[
  \{x+ie+jd: i=0,1,2,\ j=0,1\}.
\]

The cases \(e=d\) and \(e=2d\) are already excluded by Lemmas
\ref{lem:stepd} and \ref{lem:step2d}.  Thus the unresolved part of
Conjecture \ref{conj:c2} is the possibility of such grids with
\[
  e\neq d,\qquad e\neq 2d.
\]

Define \(G(A,d,N)\) to be the number of pairs \((x,e)\) such that:
\begin{itemize}
  \item \(e>0\);
  \item \(e\neq d\) and \(e\neq 2d\);
  \item \(x,x+e,x+2e\in A_d\).
\end{itemize}
Then Conjecture \ref{conj:c2} is equivalent to:
\[
  \exists d\in\calS(A)\quad G(A,d,N)=0.
\]

This is often the most useful formulation of the problem.  It says that among
popular differences, at least one column spacing \(d\) avoids all non-trivial
\(2\times3\) grids.

\section{Computational Evidence}

The repository contains several verification scripts and reports in
\texttt{runs/erdos142/}.  The core test is:

\begin{quote}
Given a 4-AP-free set \(A\subseteq\{1,\ldots,N\}\), compute all popular
differences \(d\in\calS(A)\), compute each fiber \(A_d\), and check whether at
least one popular fiber is fully 3-AP-free.
\end{quote}

\subsection{Exhaustive small-\(N\) checks}

The script \texttt{conjecture2\_verify.py} exhaustively checks all 4-AP-free
sets for small \(N\).  The project reports exhaustive verification through
\(N\leq 20\), covering approximately \(478{,}000\) 4-AP-free sets, with no
counterexample to Conjecture \ref{conj:c2}.

\subsection{Random greedy sampling}

For larger \(N\), the project uses random greedy constructions of 4-AP-free sets.
The reports state that random sampling was extended up to \(N=200000\), with no
counterexample found across approximately \(4893\) trials.

This is evidence, not proof.  Random greedy sets do not represent all possible
4-AP-free sets.  Still, the experiments are useful because they reveal which
naive strengthening attempts are false and where good fibers tend to occur.

\subsection{Observed structural patterns}

The computations show:
\begin{itemize}
  \item Conjecture \ref{conj:c2} holds in all tested examples.
  \item The largest fiber, \(d=\arg\max |A_d|\), is often not 3-AP-free.
  \item Good fibers usually occur near the lower end of the popular-fiber size
  spectrum.
  \item The fraction of bad popular fibers can be very large; reported values
  reach approximately \(0.9890\) in the largest experiments.
  \item Despite the high bad fraction, the absolute number of good popular
  fibers grows in the sampled data.
\end{itemize}

\section{A Refined Empirical Conjecture}

Let
\[
  L=\frac{M^2}{4N}
\]
be the popularity threshold and
\[
  P=\max_{d\in\calS(A)} |A_d|
\]
be the largest popular fiber size.  For \(d\in\calS(A)\), define the witness
fraction
\[
  \mathrm{wfrac}(d)=\frac{|A_d|}{P}.
\]

The computations suggest that good fibers often occur close to the popularity
threshold \(L\), and that in random greedy data
\[
  \frac{L}{P}\approx 0.20.
\]
This motivates the following stronger empirical statement.

\begin{conjecture}[Refined Conjecture 2: lower-quarter witness]
Let \(A\subseteq\{1,\ldots,N\}\) be 4-AP-free.  Then there exists
\(d\in\calS(A)\) such that \(A_d\) is fully 3-AP-free and
\[
  |A_d|\leq \frac14 \max_{d'\in\calS(A)} |A_{d'}|.
\]
\end{conjecture}

This refined conjecture is not needed for Theorem \ref{thm:conditional_vc}, but
it may be easier to attack structurally because it points toward small popular
fibers as the likely witnesses.

\section{Lean 4 Formalization}

The project includes Lean 4 files in the \texttt{lean/} directory.  The purpose
of the formalization is to check the finite combinatorial bookkeeping and to
make the gap precise.

\subsection{Formalized definitions}

The file \texttt{ArithmeticProgressions.lean} defines:
\begin{itemize}
  \item \(k\)-term arithmetic progressions as finite sets of the form
  \[
    \{a+i d:0\leq i<k\};
  \]
  \item the predicate \texttt{HasKAP A k};
  \item basic lemmas about the cardinality of such progressions;
  \item a bridge between \texttt{HasKAP A 3} and Mathlib's
  \texttt{ThreeAPFree}.
\end{itemize}

The file \texttt{RkN.lean} defines \(r_k(N)\) as a maximum over finite subsets of
\(\{0,\ldots,N-1\}\), using Lean's \texttt{Nat.findGreatest}.  It also proves the
bridge
\[
  r_3(N)=\texttt{rothNumberNat}(N),
\]
linking the custom definition to Mathlib's Roth-number definition.

\subsection{Formalized partial van Corput structure}

The file \texttt{VanCorput.lean} defines the fiber
\[
  A_d=\{x\in A:x+d\in A\}
\]
as \texttt{fiberAtDiff A d}.  It proves the Lean analogues of Lemmas
\ref{lem:stepd} and \ref{lem:step2d}.  It also proves the popular-difference
pigeonhole lemma and packages the result as a Lean theorem called
\texttt{proposition\_1}.

Thus Proposition \ref{prop:partial} is genuinely formalized.

\subsection{Declarations still using \texttt{sorry}}

The important unfinished declarations are:
\begin{itemize}
  \item \texttt{van\_corput\_inequality}, corresponding to
  \(r_4(N)\leq 2\sqrt{N r_3(N)}\);
  \item \texttt{rk\_four\_raghavan\_bound}, corresponding to the conditional
  quasi-polynomial \(r_4(N)\) bound;
  \item general Szemerédi statements for \(k\geq4\), because they are not
  available in Mathlib at the required level.
\end{itemize}

This is the correct formal status: the partial combinatorial result is proved,
while the full reduction remains blocked by Conjecture \ref{conj:c2}.

\section{What Is Actually Being Submitted}

The safe partial result is Proposition \ref{prop:partial}, together with the
isolation of Conjecture \ref{conj:c2} and the conditional consequence Theorem
\ref{thm:conditional_vc}.

The submission should not claim an unconditional new upper bound for \(r_4(N)\).
The correct claim is:

\begin{quote}
We prove that every 4-AP-free set has a popular difference fiber satisfying a
large lower bound and avoiding 3-APs of steps \(d\) and \(2d\).  We isolate the
remaining full-freeness condition as Conjecture 2.  Conditional on this
conjecture, Raghavan's \(r_3(N)\) bound implies a quasi-polynomial upper bound
for \(r_4(N)\).  Computational evidence supports the conjecture.
\end{quote}

This is a meaningful partial result because it:
\begin{enumerate}
  \item gives a clean and formally checked decomposition of the attempted van
  Corput reduction;
  \item identifies the exact obstruction rather than hiding it in an incorrect
  ``fiber is 3-AP-free'' claim;
  \item gives a concrete finite conjecture that can be tested computationally;
  \item reformulates the obstruction as a \(2\times3\) grid-avoidance problem.
\end{enumerate}

\section{Open Problems}

\begin{enumerate}
  \item Prove or disprove Conjecture \ref{conj:c2}.
  \item Prove a weaker statement still sufficient for
  \(r_4(N)\leq C\sqrt{N r_3(N)}\).
  \item Prove a positive lower bound on the number of good popular fibers, or
  show that at least one must exist by an averaging or energy argument.
  \item Establish structural restrictions on the \(2\times3\) grids counted by
  \(G(A,d,N)\).
  \item Test Conjecture \ref{conj:c2} on adversarial constructions beyond random
  greedy 4-AP-free sets.
  \item Formalize the equivalence \(G(A,d,N)=0\) iff \(A_d\) is fully 3-AP-free,
  modulo the already-proved step-\(d\) and step-\(2d\) exclusions.
\end{enumerate}

\begin{thebibliography}{9}

\bibitem{raghavan2026}
Rushil Raghavan.
\newblock Improved bounds for three-term progression-free sets.
\newblock arXiv:2603.27045, 2026.

\bibitem{greentao2017}
Ben Green and Terence Tao.
\newblock New bounds for Szemerédi's theorem, III: A polylogarithmic bound for
\(r_4(N)\).
\newblock 2017.

\bibitem{szemeredi1975}
Endre Szemerédi.
\newblock On sets of integers containing no \(k\) elements in arithmetic progression.
\newblock \emph{Acta Arithmetica}, 27:199--245, 1975.

\bibitem{ehps2024}
Christian Elsholtz, Zach Hunter, Laura Proske, and Lisa Sauermann.
\newblock Improving Behrend's construction: Sets without arithmetic progressions
in integers and over finite fields.
\newblock arXiv:2406.12290, 2024.

\bibitem{kelleymeka2023}
Zander Kelley and Raghu Meka.
\newblock Strong bounds for 3-progressions.
\newblock 2023.

\end{thebibliography}

\end{document}
